
\section*{Ventana N73: puerta de orientación}
Identidad exacta:
\[
q=u^\pi+u^e+u^\varphi,\qquad
a=u^\pi+u^e-u^\varphi.
\]
Luego:
\[
\boxed{u^\varphi=2(q-a),\qquad u^\pi+u^e=q-u^\varphi.}
\]
Así, si \(q,a\) están fijos, la ambigüedad local sólo puede vivir en el reparto de \(\pi/e\).

Auditoría \(t=5,\ldots,25\):
\[
\boxed{\text{todas las fronteras ambiguas tienen }\varphi\text{ fijo y }\pi+e\text{ fijo.}}
\]
La supervivencia coinductiva resuelve esa orientación en profundidad \(\le3\).

Novena zona:
\[
t=8,9
\]
es una puerta de orientación muy visible:
\[
\varphi\text{ fijo},\qquad \pi/e\text{ espejo móvil},\qquad \text{el futuro selecciona la orientación.}
\]
Pero no es un caso aislado: es un ejemplo emblemático de una ley más general.
